% GENERATED FILE. Edit canonical lessons, then run npm run generate:books.
% canonical-source-hash: fnv1a64:230d023e978d1701
% canonical-lessons: ZH-C166-kaishui, ZH-C166-xiawucha, ZH-W166-cha-tea, ZH-C166-hongcha, ZH-C166-chaye, ZH-C166-paocha

\chapter{Tea Time}
\label{ch:zh-tea-time}
\hlchaptermodality{166}

\begin{chapteropening}
\textbf{By the end of this chapter:} \emph{I can ask for boiled water, afternoon tea, black tea and tea leaves, and make the tea.}

Complete the last lesson of chapter 166: I can ask for boiled water, afternoon tea, black tea and tea leaves, and make the tea.
\end{chapteropening}

\section[\texorpdfstring{\zh{开水} (kāishuǐ)}{开水 (kāishuǐ)}]{\zh{开水} (kāishuǐ) — boiled water; boiling water}

\label{lesson:ZH-C166-kaishui}



\begin{quote}
Before the new one: say the Mandarin for a skirt, then the Mandarin for a scarf.
\end{quote}

\subsection*{You'll want to know: \zh{开水}}
\textbf{\zh{开水}} — \emph{kāishuǐ} — \textquotedblleft{}boiled water; boiling water\textquotedblright{}.

\begin{quote}
\textbf{\zh{中午我们吃米饭喝开水}} \emph{zhōngwǔ wǒmen chī mǐfàn hē kāishuǐ} — At midday we have rice and drink boiled water.
\end{quote}

\begin{cousinweb}
\textbf{\zh{开水}} (\emph{kāishuǐ}) is built from \zh{开} 'to open', here 'to come to the boil', and \zh{水} 'water', both inherited from Middle Chinese, \emph{khoj} and \emph{sywijX} (Guangyun, in Baxter's transcription); the boiling sense is not given there.
\end{cousinweb}

\begin{grammarlens}[title={What you've built}]
One more word for asking for what I want.
\end{grammarlens}

\subsection*{Your turn}
\begin{itemize}
\raggedright
  \item \emph{Say it:} \emph{kāishuǐ}
  \item \emph{Say it:} \emph{kāishuǐ}, once more
  \item \emph{Say it:} \emph{zhōngwǔ wǒmen chī mǐfàn hē kāishuǐ} — At midday we have rice and drink boiled water.
  \item \emph{Recall:} say the Mandarin for a skirt, then the Mandarin for a scarf, then say \emph{kāishuǐ} again
\end{itemize}

\begin{tcolorbox}[breakable,colback=teal!4,colframe=teal!35!black,title={Before you move on}]
What does \zh{开水} mean? (\textquotedblleft{}Boiled water; boiling water\textquotedblright{}.) Say it once more.
\end{tcolorbox}
\section[\texorpdfstring{\zh{下午茶} (xiàwǔchá)}{下午茶 (xiàwǔchá)}]{\zh{下午茶} (xiàwǔchá) — afternoon tea (a light meal)}

\label{lesson:ZH-C166-xiawucha}



\begin{quote}
Before the new one: say the Mandarin for a scarf, then the Mandarin for boiled water.
\end{quote}

\subsection*{You'll want to know: \zh{下午茶}}
\textbf{\zh{下午茶}} — \emph{xiàwǔchá} — \textquotedblleft{}afternoon tea (a light meal)\textquotedblright{}.

\begin{quote}
\textbf{\zh{我们三点吃下午茶}} \emph{wǒmen sān diǎn chī xiàwǔchá} — We have afternoon tea at three.
\end{quote}

The new character, in pieces: \zh{茶} — nine strokes in three pieces: the grass piece of \zh{菜} on top, the two strokes of \zh{人} under it, and at the foot a horizontal, a vertical with a hook, and two short strokes on either side of it. The next lesson puts it in your hand.

\begin{cousinweb}
\textbf{\zh{下午茶}} (\emph{xiàwǔchá}) is built from \zh{下午} 'afternoon' and \zh{茶} 'tea'. All three parts are inherited from Middle Chinese, \emph{haeX}, \emph{nguX} and \emph{drae} (Guangyun, in Baxter's transcription); the dictionary marks the tea character as a popular form. Tea with something to eat, as a light meal of its own.
\end{cousinweb}

\begin{grammarlens}[title={What you've built}]
One more word for asking for what I want.
\end{grammarlens}

\subsection*{Your turn}
\begin{itemize}
\raggedright
  \item \emph{Say it:} \emph{xiàwǔchá}
  \item \emph{Say it:} \emph{xiàwǔchá}, once more
  \item \emph{Say it:} \emph{wǒmen sān diǎn chī xiàwǔchá} — We have afternoon tea at three.
  \item \emph{Recall:} say the Mandarin for a scarf, then the Mandarin for boiled water, then say \emph{xiàwǔchá} again
\end{itemize}

\begin{tcolorbox}[breakable,colback=teal!4,colframe=teal!35!black,title={Before you move on}]
What does \zh{下午茶} mean? (\textquotedblleft{}Afternoon tea (a light meal)\textquotedblright{}.) Say it once more.
\end{tcolorbox}
\section[\texorpdfstring{\zh{茶} (chá)}{茶 (chá)}]{\zh{茶} — grass, a roof, a small tree}

\label{lesson:ZH-W166-cha-tea}



\begin{quote}
The word first, then the character inside it.

\begin{itemize}
\raggedright
  \item \emph{Recall:} say \emph{xiàwǔchá} — \textbf{R1}, one lesson back, before it decays
\end{itemize}

\begin{quote}
\zh{菜} \zh{人}
\end{quote}

\emph{Read it:} each one aloud, with nothing said first — \emph{cài}, \emph{rén}
\end{quote}

\subsection*{Script — grass, a roof, a small tree}
\hlblockfigure{\detokenize{figures/ZH-W166-cha-tea-filmstrip.pdf}}{How \zh{茶} is written, stroke by stroke}

\begin{quote}
\zh{茶}
\end{quote}

\textbf{Three pieces, top to bottom.} Count the pieces, not the strokes.

\textbf{The grass piece, three strokes,} as in \zh{菜}: a horizontal, then two short verticals through it.

\textbf{Then \zh{人}, two strokes,} spread wide like a roof.

\textbf{At the foot, four strokes:} a horizontal, a vertical with a hook, then a short stroke falling to the left and a dot falling to the right, one on each side.

3 + 2 + 4 = nine.

\subsection*{Your turn}
\begin{itemize}
\raggedright
  \item \emph{Write it:} \zh{茶}, all nine, with the model in front of you
  \item \emph{Cover:} write it again from memory, and count the pieces, not the strokes
\end{itemize}

\begin{tcolorbox}[breakable,colback=teal!4,colframe=teal!35!black,title={Before you move on}]
How many strokes in \zh{茶}? (\textbf{Nine}.) What sits on top? (\textbf{The grass piece of \zh{菜}}.) What comes under it? (\textbf{\zh{人}, spread like a roof}.)

Source: \href{https://raw.githubusercontent.com/chanind/hanzi-writer-data/68d10a4b21150cae5e1ebbd223eed289cf32d90c/data/\%E8\%8C\%B6.json}{Hanzi Writer Data \zh{茶}, pinned stroke paths and medians}.
\end{tcolorbox}
\section[hóngchá]{hóngchá — black tea}

\label{lesson:ZH-C166-hongcha}



\begin{quote}
Before the new one: say the Mandarin for boiled water, then the Mandarin for afternoon tea.
\end{quote}

\subsection*{You'll want to know: hóngchá}
\emph{hóngchá} — \textquotedblleft{}black tea\textquotedblright{}.

\begin{quote}
\emph{nǐ yào hóngchá háishi lǜchá} — Black tea or green tea?
\end{quote}

Its characters come later; for now it is a word for the ear.

English calls this tea black; Mandarin calls it red.

\begin{cousinweb}
\emph{hóngchá} is built from \emph{hóng} 'red' and \zh{茶} 'tea', both inherited from Middle Chinese, \emph{huwng} and \emph{drae} (Guangyun, in Baxter's transcription).
\end{cousinweb}

\begin{grammarlens}[title={What you've built}]
One more word for asking for what I want.
\end{grammarlens}

\subsection*{Your turn}
\begin{itemize}
\raggedright
  \item \emph{Say it:} \emph{hóngchá}
  \item \emph{Say it:} \emph{hóngchá}, once more
  \item \emph{Say it:} \emph{nǐ yào hóngchá háishi lǜchá} — Black tea or green tea?
  \item \emph{Recall:} say the Mandarin for boiled water, then the Mandarin for afternoon tea, then say \emph{hóngchá} again
\end{itemize}

\begin{tcolorbox}[breakable,colback=teal!4,colframe=teal!35!black,title={Before you move on}]
What does hóngchá mean? (\textquotedblleft{}Black tea\textquotedblright{}.) Say it once more.
\end{tcolorbox}
\section[cháyè]{cháyè — tea leaves}

\label{lesson:ZH-C166-chaye}



\begin{quote}
Before the new one: say the Mandarin for afternoon tea, then the Mandarin for black tea.
\end{quote}

\subsection*{You'll want to know: cháyè}
\emph{cháyè} — \textquotedblleft{}tea leaves\textquotedblright{}.

\begin{quote}
\emph{cháyè zài nǎr} — Where are the tea leaves?
\end{quote}

Its characters come later; for now it is a word for the ear.

\begin{cousinweb}
\emph{cháyè} is built from \zh{茶} 'tea' and \emph{yè} 'leaf', both inherited from Middle Chinese, \emph{drae} and \emph{yep} (Guangyun, in Baxter's transcription), which glosses the second 'the leaves of a branch'.
\end{cousinweb}

\begin{grammarlens}[title={What you've built}]
One more word for asking for what I want.
\end{grammarlens}

\subsection*{Your turn}
\begin{itemize}
\raggedright
  \item \emph{Say it:} \emph{cháyè}
  \item \emph{Say it:} \emph{cháyè}, once more
  \item \emph{Say it:} \emph{cháyè zài nǎr} — Where are the tea leaves?
  \item \emph{Recall:} say the Mandarin for afternoon tea, then the Mandarin for black tea, then say \emph{cháyè} again
\end{itemize}

\begin{tcolorbox}[breakable,colback=teal!4,colframe=teal!35!black,title={Before you move on}]
What does cháyè mean? (\textquotedblleft{}Tea leaves\textquotedblright{}.) Say it once more.
\end{tcolorbox}
\section[pàochá]{pàochá — to make tea (to steep the leaves)}

\label{lesson:ZH-C166-paocha}



\begin{quote}
Before the new one: say the Mandarin for black tea, then the Mandarin for tea leaves.
\end{quote}

\subsection*{You'll want to know: pàochá}
\emph{pàochá} — \textquotedblleft{}to make tea (to steep the leaves)\textquotedblright{}.

\begin{quote}
\emph{wǒ yòng kāishuǐ pàochá} — I make the tea with boiling water.
\end{quote}

Its characters come later; for now it is a word for the ear.

\begin{cousinweb}
\emph{pàochá} is built from \emph{pào} 'to soak' and \zh{茶} 'tea', both inherited from Middle Chinese, \emph{phaew} and \emph{drae} (Guangyun, in Baxter's transcription). The dictionary glosses the first 'bubbles on water'; soaking is not given there, and \emph{phaew} would regularly give a first tone, so how \emph{pào} came about is not settled.
\end{cousinweb}

\begin{grammarlens}[title={What you've built}]
One more word for asking for what I want.
\end{grammarlens}

\subsection*{Your turn}
\begin{itemize}
\raggedright
  \item \emph{Say it:} \emph{pàochá}
  \item \emph{Say it:} \emph{pàochá}, once more
  \item \emph{Say it:} \emph{wǒ yòng kāishuǐ pàochá} — I make the tea with boiling water.
  \item \emph{Recall:} say the Mandarin for black tea, then the Mandarin for tea leaves, then say \emph{pàochá} again
\end{itemize}

\begin{tcolorbox}[breakable,colback=teal!4,colframe=teal!35!black,title={Before you move on}]
What does pàochá mean? (\textquotedblleft{}To make tea (to steep the leaves)\textquotedblright{}.) Say it once more.
\end{tcolorbox}
